\ifdefined\gkver\else\def\gkver{student}\fi
\documentclass[\gkver]{gaokaozhenti}
\begin{document}

\chapter{2012年天津理综}
%% paperType: 理综物理部分

\begin{enumerate}
\item
%% number: 1
%% paperName: 2012年天津理综第 1 题 6 分
%% typeId: 1
%% type: 单选题
%% chapter: 原子和原子核
%% point: 玻尔模型
%% method: 无
%% score: 6
%% degree: 650
%% duplicateId: 0
%% body:
下列说法正确的是\xzanswer{B}
\fourchoices[answer=B]
{采用物理或化学方法可以有效地改变放射性元素的半衰期}
{由波尔理论知道氢原子从激发态跃迁到基态时会放出光子}
{从高空对地面进行遥感摄影是利用紫外线良好的穿透能力}
{原子核所含核子单独存在时的总质量小于该原子核的质量}
%% answer: B
%% memo:
\memoanswer{
半衰期是原子的物理属性，不能采用物理或化学方法改变；高空遥感是用红外线的；由于核子结合为原子核时能量增加必然存在质量亏损；氢原子从高能量的激发态跃迁到低能量的基态时放出能量，所以放出光子。答案 B。
}

\item
%% number: 2
%% paperName: 2012年天津理综第 2 题 6 分
%% typeId: 1
%% type: 单选题
%% chapter: 磁场
%% point: 安培力
%% method: 无
%% score: 6
%% degree: 650
%% duplicateId: 0
%% body:
如图所示，金属棒 MN 两端由等长的轻质细线水平悬挂，处于竖直向上的匀强磁场中，棒中通以由 M 向 N 的电流，平衡时两悬线与竖直方向夹角均为 $\theta$，如果仅改变下列某一个条件，$\theta$ 角的相应变化情况是\xzanswer{A}
\onepicture[h=2.6cm, align=right]{figs/02.png}
\fourchoices[answer=A]
{棒中的电流变大，$\theta$ 角变大}
{两悬线等长变短，$\theta$ 角变小}
{金属棒质量变大，$\theta$ 角变大}
{磁感应强度变大，$\theta$ 角变小}
%% answer: A
%% memo:
\memoanswer{
水平的直线电流在竖直磁场中受到水平的安培力而偏转，与竖直方向形成夹角，此时它受拉力、重力和安培力而达到平衡，根据平衡条件有 $\tan\theta=\frac{F_{\text{安}}}{mg}=\frac{BIL}{mg}$，所以棒中的电流增大，$\theta$ 角变大；两悬线变短，不影响平衡状态，$\theta$ 角不变；金属棒质量变大，$\theta$ 角变小；磁感应强度变大，$\theta$ 角变大。答案 A。
}

\item
%% number: 3
%% paperName: 2012年天津理综第 3 题 6 分
%% typeId: 1
%% type: 单选题
%% chapter: 曲线运动
%% point: 人造地球卫星
%% method: 无
%% score: 6
%% degree: 450
%% duplicateId: 0
%% body:
一人造地球卫星绕地球做匀速圆周运动，假如该卫星变轨后做匀速圆周运动，动能减小为原来的 1/4，不考虑卫星质量的变化，则变轨前后卫星的\xzanswer{C}
\fourchoices[answer=C]
{向心加速度大小之比为 4$:$1}
{角速度大小之比为 2$:$1}
{周期之比为 1$:$8}
{轨道半径之比为 1$:$2}
%% answer: C
%% memo:
\memoanswer{
根据向心加速度表达式 $a = \frac{v^2}{R}$ 知在动能减小时势能增大，地球卫星的轨道半径增大，则向心加速度之比大于 4；根据万有引力和牛顿第二定律有 $m \frac{v^2}{R} = G \frac{Mm}{R^2}$，化简为 $Rv^2 = GM$，知在动能减小、速度减小时轨道半径增大到原来的 4 倍；同理有 $m \left( \frac{2\pi}{T} \right)^2 R = G \frac{Mm}{R^2}$，化简为 $\frac{R^3}{T^2} = \frac{GM}{4\pi^2}$，则周期的平方增大到 8 倍；根据角速度关系式 $\omega = \frac{2\pi}{T}$，角速度减小为 $\frac{1}{8}$。答案 C。
}

\item
%% number: 4
%% paperName: 2012年天津理综第 4 题 6 分
%% typeId: 1
%% type: 单选题
%% chapter: 交流电
%% point: 变压器
%% method: 无
%% score: 6
%% degree: 450
%% duplicateId: 0
%% body:
通过一理想变压器，经同一线路输送相同电功率 $P$，原线圈的电压 $U$ 保持不变，输电线路的总电阻为 $R$。当副线圈与原线圈的匝数比为 $k$ 时，线路损耗的电功率为 $P_{1}$，若将副线圈与原线圈的匝数比提高到 $nk$，线路损耗的电功率为 $P_{2}$，则 $P_{1}$ 和 $\frac{P_{2}}{P_{1}}$ 分别为\xzanswer{D}
\fourchoices[answer=D]
{$\frac{PR}{kU}$，$\frac{1}{n}$}
{$\left(\frac{P}{kU}\right)^{2}R$，$\frac{1}{n}$}
{$\frac{PR}{kU}$，$\frac{1}{n^{2}}$}
{$\left(\frac{P}{kU}\right)^{2}R$，$\frac{1}{n^{2}}$}
%% answer: D
%% memo:
\memoanswer{
根据理想变压器的变压比有 $k=\frac{U_{1}}{U}$，$nk=\frac{U_{2}}{U}$。线路的输送功率不变，有 $P = U_{1} I_{1} = U_{2} I_{2}$。根据焦耳定律有 $P_{1}=I_{1}^{2}R=\left(\frac{P}{U_{1}}\right)^{2}R=\left(\frac{P}{kU}\right)^{2}R$，$P_{2}=I_{2}^{2}R$，$\frac{P_{2}}{P_{1}}=\frac{I_{2}^{2}}{I_{1}^{2}}=\frac{U_{1}^{2}}{U_{2}^{2}}=\frac{1}{n^{2}}$。答案 D。
}

\item
%% number: 5
%% paperName: 2012年天津理综第 5 题 6 分
%% typeId: 1
%% type: 单选题
%% chapter: 电场
%% point: 等势面
%% method: 无
%% score: 6
%% degree: 450
%% duplicateId: 0
%% body:
两个固定的等量异号点电荷所产生电场的等势面如图中虚线所示，一带负电的粒子以某一速度从图中 A 点沿图示方向进入电场在纸面内飞行，最后离开电场，粒子只受静电力作用，则粒子在电场中\xzanswer{C}
\onepicture[h=2.6cm, align=right]{figs/05.png}
\fourchoices[answer=C]
{做直线运动，电势能先变小后变大}
{做直线运动，电势能先变大后变小}
{做曲线运动，电势能先变小后变大}
{做曲线运动，电势能先变大后变小}
%% answer: C
%% memo:
\memoanswer{
两个固定的等量异号点电荷产生电场的等势面与电场线垂直，且沿电场线电势减小，所以等量异号点电荷和它们间的直线电场线如图所示。当带负电的粒子进入电场受到与其速度垂直的电场力而偏转做曲线运动。电子在斜向左的电场力作用下向左偏转，电场力做正功，粒子的动能增加，电势能减小。答案 C。
}

\item
%% number: 6
%% paperName: 2012年天津理综第 6 题 6 分
%% typeId: 2
%% type: 多选题
%% chapter: 光学
%% point: 光的折射
%% method: 无
%% score: 6
%% degree: 650
%% duplicateId: 0
%% body:
半圆形玻璃砖横截面如图，AB 为直径，O 点为圆心。在该截面内有 a 、b 两束单色可见光从空气垂直于 AB 射入玻璃砖，两入射点到 O 的距离相等。两束光在半圆边界上反射和折射的情况如图所示：则 a、b 两束光，\xzanswer{ACD}
\onepicture[h=3.0cm, align=right]{figs/06.png}
\fourchoices[answer=ACD]
{在同种均匀介质中传播，a 光的传播速度较大}
{以相同的入射角从空气斜射入水中，b 光的折射角大}
{若 a 光照射某金属表面能发生光电效应，b 光也一定能}
{分别通过同一双缝干涉装置，a 光的相邻亮条纹间距大}
%% answer: ACD
%% memo:
\memoanswer{
当光由光密介质——玻璃进入光疏介质——空气时发生折射或全反射，b 发生全反射说明 b 的入射角大于或等于临界角，a 发生折射说明 a 的入射角小于临界角，比较可知在玻璃中 a 的临界角大于 b 的临界角；根据临界角定义有 $\sin C = \frac{1}{n}$，玻璃对 a 的折射率小；根据 $n = \frac{c}{v}$，在玻璃中 a 光的速度大，A 正确；通过色散现象分析比较，a 的折射率小，a 光的频率小、波长大；双缝干涉相邻亮条纹间距大小与波长成正比，a 光的相邻亮条纹间距大，D 正确；发生光电效应的条件是入射光的频率大于金属的极限频率，频率小的 a 光能发生光电效应，则频率大的 b 光一定能，C 正确；根据折射定律 $n = \frac{\sin i}{\sin r}$，在入射角 i 相同时 b 的折射率大则折射角 r 小，B 错误。答案 ACD。
}

\item
%% number: 7
%% paperName: 2012年天津理综第 7 题 6 分
%% typeId: 2
%% type: 多选题
%% chapter: 机械波
%% point: 波的图象
%% method: 无
%% score: 6
%% degree: 450
%% duplicateId: 0
%% body:
沿 $x$ 轴正向传播的一列简谐横波在 $t=$0 时刻的波形如图所示，M 为介质中的一个质点，该波的传播速度为 $40\Ums$，则 $t=\frac{1}{40}\Us$ 时\xzanswer{CD}
\onepicture[h=2.4cm, align=right]{figs/07.png}
\fourchoices[answer=CD]
{质点 M 对平衡位置的位移一定为负值}
{质点 M 的速度方向与对平衡位置的位移方向相同}
{质点 M 的加速度方向与速度方向一定相同}
{质点 M 的加速度方向与对平衡位置的位移方向相反}
%% answer: CD
%% memo:
\memoanswer{
由波的图像知波长是 $4\Um$，所以波源的振动周期是 $T = \frac{\lambda}{v} = \frac{4}{40}\Us = \frac{1}{10}\Us$，则 $t = \frac{T}{4}$。波是沿 $x$ 轴正方向传播，可判断质点 M 在 0 时刻沿 $y$ 轴正方向振动，$t$ 时刻在 $x$ 轴上方运动，位移为正，速度为负，加速度为负，AB 错误，CD 正确。答案 CD。
}

\item
%% number: 8
%% paperName: 2012年天津理综第 8 题 6 分
%% typeId: 2
%% type: 多选题
%% chapter: 运动和力
%% point: 牛顿定律中的图象
%% method: 无
%% score: 6
%% degree: 450
%% duplicateId: 0
%% body:
如图甲所示，静止在水平地面的物块 A，受到水平向右的拉力 $F$ 的作用，$F$ 与时间 $t$ 的关系如图乙所示，设物块与地面的静摩擦力最大值 $f_{m}$ 与滑动摩擦力大小相等，则\xzanswer{BD}
\onepicture[h=2.2cm, align=right]{figs/08.png}
\fourchoices[answer=BD]
{0\textasciitilde{}$t_{1}$ 时间内 $F$ 的功率逐渐增大}
{$t_{2}$ 时刻物块 A 的加速度最大}
{$t_{2}$ 时刻后物块 A 做反向运动}
{$t_{3}$ 时刻物块 A 的动能最大}
%% answer: BD
%% memo:
\memoanswer{
由 $F$ 与 $t$ 的关系图像，$0\sim t_1$ 拉力小于最大静摩擦力，物块静止，$F$ 的功率为 0，A 错误；在 $t_1\sim t_2$ 阶段拉力大于最大静摩擦力，物块做加速度增大的加速运动，在 $t_2\sim t_3$ 阶段拉力大于最大静摩擦力，物块做加速度减小的加速运动，在 $t_2$ 时刻加速度最大，B 正确，C 错误；在 $t_1\sim t_3$ 物块一直做加速运动，在 $t_3\sim t_4$ 拉力小于最大静摩擦力，物块开始减速，在 $t_3$ 时刻速度最大，动能最大，D 正确。答案 BD。
}

\item
%% number: 9
%% paperName: 2012年天津理综第 9 题 4 分
%% typeId: 3
%% type: 填空题
%% chapter: 运动和力
%% point: 动量定理
%% method: 无
%% score: 4
%% degree: 650
%% duplicateId: 0
%% body:
质量为 $0.2\Ukg$ 的小球竖直向下以 $6\Ums$ 的速度落至水平地面，再以 $4\Ums$ 的速度反向弹回，取竖直向上为正方向，则小球与地面碰撞前后的动量变化为\_\_\_\_\_\_\_\_$\Ukgms$。若小球与地面的作用时间为 $0.2\Us$，则小球受到地面的平均作用力大小为\_\_\_\_\_\_\_\_$\UN$（$g=10\Umsq$）。
%% answer: 
\jdanswer{
2，12
}
%% memo:
\memoanswer{
取竖直向上为正方向，则初动量为负、末动量为正，动量变化为 $\Delta p = p' - p = 4 \times 0.2 - (-6 \times 0.2)\Ukgms = 2\Ukgms$，$F = \frac{\Delta p}{t} + mg = \frac{2}{0.2}\UN + 0.2 \times 10\UN = 12\UN$。
}

\item
%% number: 9
%% paperName: 2012年天津理综第 9 题 6 分
%% typeId: 4
%% type: 实验题
%% chapter: 机械振动
%% point: 单摆测重力加速度
%% method: 无
%% score: 6
%% degree: 450
%% duplicateId: 0
%% body:
某同学用实验的方法探究影响单摆周期的因素。
\begin{enumerate}
	\item
	他组装单摆时，在摆线上端的悬点处，用一块开有狭缝的橡皮夹牢摆线，再用铁架台的铁夹将橡皮夹紧，如图所示。这样做的目的是\_\_\_\_\_\_\_\_（填字母代号）。

	（A）保证摆动过程中摆长不变\quad（B）可使周期测量得更加准确

	（C）需要改变摆长时便于调节\quad（D）保证摆球在同一竖直平面内摆动
	\onepicture[h=2.6cm, align=right]{figs/09a.png}

	\item
	他组装好单摆后在摆球自然悬垂的情况下，用毫米刻度尺从悬点量到摆球的最低端的长度 $L=0.9990\Um$，再用游标卡尺测量摆球直径，结果如图所示，则摆球的直径为\_\_\_\_\_\_\_\_$\Umm$，单摆摆长为\_\_\_\_\_\_\_\_$\Um$。
	\onepicture[h=2.0cm, align=right]{figs/09b.png}

	\item
	下列摆动图像真实地描述了对摆长约为 $1\Um$ 的单摆进行周期测量的四种操作过程，图中横坐标原点表示计时开始，A、B、C 均为 30 次全振动图象，已知 $\sin 5^{\circ}=0.087$，$\sin 15^{\circ}=0.26$，这四种操作过程合乎实验要求且误差最小的是\_\_\_\_\_\_\_\_（填字母代号）。
	\onepicture[w=13cm]{figs/09c.png}
\end{enumerate}
%% answer: 
\jdanswer{
\begin{enumerate}
	\item
	AC

	\item
	12.0，0.9930

	\item
	A

\end{enumerate}
}
%% memo:
\memoanswer{
① 用一块开有狭缝的橡皮夹牢摆线，增加了线与悬挂处的摩擦，保证摆长不变；改变摆长时用力拉不会将摆线拉断，方便调节摆长；AC 正确。

② 用 10 分度游标卡尺，主尺读数为 12 mm，游标尺读数为 0，则摆球的直径为 $12\Umm + 0 \times 0.1\Umm = 12.0\Umm$。单摆摆长为 $l = L - \frac{D}{2} = 0.9990\Um - \frac{0.012}{2}\Um = 0.9930\Um$。

③ 单摆的振动在摆角小于 5 度才能看作简谐振动，在测量周期时计时起点应该选择在平衡位置（速度大误差小）。根据摆角估算振幅 $A_1 = l \sin 5^{\circ} = 0.993 \times 0.087\Um = 0.086\Um = 8.6\Ucm$，A、B 振幅合理。$A_2 = l \sin 15^{\circ} = 0.993 \times 0.26\Um = 0.25\Um = 25\Ucm$，C、D 振幅不合理，错误。A 中振动图像的计时起点在平衡位置是合理的，B 中振动图像的计时起点在正的最大位置是不合理的。答案 A。
}

\item
%% number: 9
%% paperName: 2012年天津理综第 9 题 8 分
%% typeId: 4
%% type: 实验题
%% chapter: 电路
%% point: 电路实验
%% method: 无
%% score: 8
%% degree: 450
%% duplicateId: 0
%% body:
某同学在进行扩大电流表量程的实验时，需要知道电流表的满偏电流和内阻。他设计了一个用标准电流表 G$_{1}$ 来校对待测电流表 G$_{2}$ 的满偏电流和测定 G$_{2}$ 内阻的电路，如图所示。已知 G$_{1}$ 的量程略大于 G$_{2}$ 的量程，图中 $R_{1}$ 为滑动变阻器，$R_{2}$ 为电阻箱。该同学顺利完成了这个实验。
\begin{enumerate}
	\item
	实验过程包含以下步骤，其合理的顺序依次为\_\_\_\_\_\_\_\_\_\_\_\_\_\_\_\_\_\_（填步骤的字母代号）。
	\onepicture[h=2.6cm, align=right]{figs/09.png}

	（A）合上开关 $S_{2}$

	（B）分别将 $R_{1}$ 和 $R_{2}$ 的阻值调至最大

	（C）记下 $R_{2}$ 的最终读数

	（D）反复调节 $R_{1}$ 和 $R_{2}$ 的阻值，使 $G_{1}$ 的示数仍为 $I_{1}$，使 $G_{2}$ 的指针偏转到满刻度的一半，此时 $R_{2}$ 的最终读数为 r

	（E）合上开关 $S_{1}$

	（F）调节 $R_{1}$ 使 $G_{2}$ 的指针偏转到满刻度，此时 $G_{1}$ 的示数为 $I_{1}$，记下此时 $G_{1}$ 的示数

	\item
	仅从实验设计原理上看，用上述方法得到的 G$_{2}$ 内阻的测量值与真实值相比\_\_\_\_\_\_\_\_（填“偏大”、“偏小”或“相等”）。

	\item
	若要将 G$_{2}$ 的量程扩大为 $I$，并结合前述实验过程中测量得结果，写出须在 G$_{2}$ 上并联的分流电阻 $R_{S}$ 的表达式，$R_{S}=$\_\_\_\_\_\_\_\_。
\end{enumerate}
%% answer: 
\jdanswer{
\begin{enumerate}
	\item
	BEFADC

	\item
	相等

	\item
	$\frac{I_{1}}{I - I_{1}}r$

\end{enumerate}
}
%% memo:
\memoanswer{
① 在组装实验时正确操作是将所有开关断开，将滑动变阻器和电阻箱的电阻调到最大，然后闭合电源控制开关 $S_{1}$，就可以校对待测表 $G_{2}$。由于标准表和待测表串联流过的电流相等，这时就可以调节 $R_{1}$，由小到大改变电流对待测表的刻度逐个校对，直到待测表满偏。校对完成后闭合开关 $S_{2}$ 构成半偏法测量表头 $G_{2}$ 内阻的电路，在没有 $G_{1}$ 的情况下只调节 $R_{2}$ 使待测表半偏，则表头内阻等于 $R_{2}$ 此时的读数，这种做法会造成 $G_{1}$ 的示数增大（原理上是不变），误差大。由于有 $G_{1}$ 可以方便观察干路电流的变化，所以操作为闭合开关 $S_{2}$，反复调节 $R_{1}$ 和 $R_{2}$，使标准表达到待测表满偏时的电流值，同时使待测表半偏，则表头内阻等于 $R_{2}$ 此时的读数。正确顺序为 B、E、F、A、D、C。

② 由于 $G_{1}$ 测量了干路电流，$G_{2}$ 测量了支路电流，另一支路电流是两表示数之差，然后根据并联电路电流、电阻成反比准确计算。答案：相等。

③ 扩大电流表的量程需要并联一个电阻，根据并联关系有 $rI_{1}=(I-I_{1})R_{S}$，则 $R_{S}=\frac{I_{1}}{I-I_{1}}r$。
}

\item
%% number: 10
%% paperName: 2012年天津理综第 10 题 16 分
%% typeId: 6
%% type: 计算题
%% chapter: 运动和力
%% point: 动量守恒定律
%% method: 无
%% score: 16
%% degree: 450
%% duplicateId: 0
%% body:
如图所示，水平地面上固定有高为 $h$ 的平台，台面上有固定的光滑坡道，坡道顶端距台面高也为 $h$，坡道底端与台面相切。小球 A 从坡道顶端由静止开始滑下，到达水平光滑的台面后与静止在台面上的小球 B 发生碰撞，并粘连在一起，共同沿台面滑行并从台面边缘飞出，落地点与飞出点的水平距离恰好为台高的一半。两球均可视为质点，忽略空气阻力，重力加速度为 $g$。求：
\begin{enumerate}
	\item
	小球 A 刚滑至水平台面的速度 $v_{A}$；
	\item
	A、B 两球的质量之比 $m_{A}:m_{B}$。
\end{enumerate}
\onepicture[h=2.8cm, align=right]{figs/10.png}
%% answer: 
\jdanswer{
\begin{enumerate}
	\item
	$\sqrt{2gh}$

	\item
	1:3

\end{enumerate}
}
%% memo:
\memoanswer{
（1）小球 A 在坡道上只有重力做功，机械能守恒，有 $\frac{1}{2}m_{A}v_{A}^{2}=m_{A}gh$，解得 $v_{A}=\sqrt{2gh}$。

（2）小球 A、B 在光滑台面上发生碰撞粘在一起，速度为 v，根据系统动量守恒得 $m_{A}v_{A}=(m_{A}+m_{B})v$。离开平台后做平抛运动，在竖直方向有 $\frac{1}{2}gt^{2}=h$，在水平方向有 $\frac{1}{2}h = vt$。联立以上各式化简得 $m_{A}:m_{B}=1:3$。
}

\item
%% number: 11
%% paperName: 2012年天津理综第 11 题 18 分
%% typeId: 6
%% type: 计算题
%% chapter: 电磁感应
%% point: 电磁感应能量分析
%% method: 无
%% score: 18
%% degree: 450
%% duplicateId: 0
%% body:
如图所示，一对光滑的平行金属导轨固定在同一水平面内，导轨间距 $L=0.5\Um$，左端接有阻值 $R=0.3\UO$ 的电阻，一质量 $m=0.1\Ukg$、电阻 $r=0.1\UO$ 的金属棒 MN 放置在导轨上，整个装置置于竖直向上的匀强磁场中，磁场的磁感应强度 $B=0.4\UT$。棒在水平向右的外力作用下，由静止开始以 $a=2\Umsq$ 的加速度做匀加速运动，当棒的位移 $x=9\Um$ 时撤去外力，棒继续运动一段距离后停下来，已知撤去外力前后回路中产生的焦耳热之比 $Q_{1}:Q_{2}=$2$:$1。导轨足够长且电阻不计，棒在运动过程中始终与导轨垂直且两端与导轨保持良好接触。求：
\begin{enumerate}
	\item
	棒在匀加速运动过程中，通过电阻 R 的电荷量 $q$；
	\item
	撤去外力后回路中产生的焦耳热 $Q_{2}$；
	\item
	外力做的功 $W_{F}$。
\end{enumerate}
\onepicture[h=2.6cm, align=right]{figs/11.png}
%% answer: 
\jdanswer{
\begin{enumerate}
	\item
	$4.5\UC$

	\item
	$1.8\UJ$

	\item
	$5.4\UJ$

\end{enumerate}
}
%% memo:
\memoanswer{
（1）棒匀加速运动所用时间为 t，有 $\frac{1}{2}at^{2}=x$，得 $t=\sqrt{\frac{2x}{a}}=\sqrt{\frac{2\times9}{2}}\Us=3\Us$。根据法拉第电磁感应定律和闭合电路的欧姆定律，电路中产生的平均电流为 $\overline{I}=\frac{\overline{E}}{R+r}=\frac{\Delta\Phi}{t(r+R)}=\frac{Blx}{t(r+R)}=\frac{0.4\times0.5\times9}{3\times(0.3+0.1)}\UA=1.5\UA$。根据电流定义式有 $q=\overline{I}t=1.5\times3\UC=4.5\UC$。

（2）撤去外力前棒做匀加速运动，根据速度公式末速为 $v=at=2\times3\Ums=6\Ums$。撤去外力后棒在安培力作用下做减速运动，安培力做负功先将棒的动能转化为电能，再通过电流做功将电能转化为内能，所以焦耳热等于棒的动能减少量。有 $Q_{2}=\Delta E_{k}=\frac{1}{2}mv^{2}=\frac{1}{2}\times0.1\times6^{2}\UJ=1.8\UJ$。

（3）根据题意，在撤去外力前的焦耳热为 $Q_{1}=2Q_{2}=3.6\UJ$。撤去外力前拉力做正功、安培力做负功（其大小等于焦耳热 $Q_{1}$）、重力不做功，共同使棒的动能增大，根据动能定理有 $\Delta E_{k}=W_{F}-Q_{1}$。则 $W_{F}=Q_{1}+\Delta E_{k}=3.6\UJ+1.8\UJ=5.4\UJ$。
}

\item
%% number: 12
%% paperName: 2012年天津理综第 12 题 20 分
%% typeId: 6
%% type: 计算题
%% chapter: 磁场
%% point: 带电粒子在组合场中的运动
%% method: 无
%% score: 20
%% degree: 450
%% duplicateId: 0
%% body:
对铀 235 的进一步研究在核能的开发和利用中具有重要意义。如图所示，质量为 $m$、电荷量为 $q$ 的铀 235 离子，从容器 A 下方的小孔 S$_{1}$ 不断飘入加速电场，其初速度可视为零，然后经过小孔 S$_{2}$ 垂直于磁场方向进入磁感应强度为 $B$ 的匀强磁场中，做半径为 R 的匀速圆周运动，离子行进半个圆周后离开磁场并被收集，离开磁场时离子束的等效电流为 $I$。不考虑离子重力及离子间的相互作用。
\begin{enumerate}
	\item
	求加速电场的电压 $U$；
	\item
	求出在离子被收集的过程中任意时间 $t$ 内收集到离子的质量 $M$；
	\item
	实际上加速电压的大小会在 $U\pm\Delta U$ 范围内微小变化。若容器 A 中有电荷量相同的铀 235 和铀 238 两种离子，如前述情况它们经电场加速后进入磁场中发生分离，为使这两种离子在磁场中运动的轨迹不发生交叠，$\frac{\Delta U}{U}$ 应小于多少？（结果用百分数表示，保留两位有效数字）
\end{enumerate}
\onepicture[h=2.8cm, align=right]{figs/12.png}
%% answer: 
\jdanswer{
\begin{enumerate}
	\item
	$\frac{qB^{2}R^{2}}{2m}$

	\item
	$\frac{mIt}{q}$

	\item
	0.63\%

\end{enumerate}
}
%% memo:
\memoanswer{
（1）铀粒子在电场中加速到速度 v，根据动能定理有 $\frac{1}{2}mv^{2}=qU$。进入磁场后在洛伦兹力作用下做圆周运动，根据牛顿第二定律有 $\frac{mv^{2}}{R}=qvB$。由以上两式化简得 $U=\frac{qB^{2}R^{2}}{2m}$。

（2）在时间 t 内收集到的粒子个数为 N，粒子总电荷量为 Q，则 $Q = It$，$N = \frac{Q}{q}$，$M = Nm$，解得 $M = \frac{mIt}{q}$。

（3）两种粒子在磁场中运动的轨迹不发生交叠，即不要重合，可得半径为 $R = \frac{1}{B} \sqrt{\frac{2mU}{q}}$。由此可知质量小的铀 235 在电压最大时的半径存在最大值 $R_{\max} = \frac{1}{B} \sqrt{\frac{2m(U + \Delta U)}{q}}$；质量大的铀 238 质量 $m'$ 在电压最小时的半径存在最小值 $R_{\min} = \frac{1}{B} \sqrt{\frac{2m'(U - \Delta U)}{q}}$。所以两种粒子在磁场中运动的轨迹不发生交叠的条件为 $\frac{1}{B} \sqrt{\frac{2m(U + \Delta U)}{q}} < \frac{1}{B} \sqrt{\frac{2m'(U - \Delta U)}{q}}$。化简得 $\frac{\Delta U}{U} < \frac{m' - m}{m' + m} = \frac{238u - 235u}{238u + 235u} = \frac{3}{473} = 0.63\%$。
}

\end{enumerate}

\end{document}
